\documentclass[11pt]{article}
\usepackage{fontspec,amsmath,amssymb,amsthm,hyperref}
\usepackage[margin=1in]{geometry}
\setlength{\emergencystretch}{3em}
\title{Orthogonal masks and the exact removed-noise cost}
\author{LangevinLean mathematical exposition}
\date{}
\begin{document}\maketitle
\noindent Source module: \texttt{ProximalLangevin/MYPGDNoise.lean}. Compile this proof document with XeLaTeX or LuaLaTeX. The accompanying HTML pairs the mathematical argument with the exact Lean statements and records the scope of the formalization.


Let \(P_\theta\) retain the common parameter block of the scaled particle state, and \(P_x=I-P_\theta\) retain the latent blocks. For a standard Gaussian vector \(\xi\), the removed parameter noise has mean zero and second moment \(d_\theta\).


\section{Define the actual complementary coordinate projections}


Set \(P_\theta(\theta,x)=(\theta,0)\) and \(P_x(\theta,x)=(0,x)\). Coordinate evaluation proves the four block identities by unfolding the definitions. These are bounded linear maps, so masking an \(L^2\) innovation remains an actual \(L^2\) random variable.




Denote the two maps by \(P_\theta\) and \(P_x=I-P_\theta\). On a scaled particle state, \(P_\theta(\theta,z_1,\ldots,z_N)=(\theta,0,\ldots,0)\), while \(P_x(\theta,z_1,\ldots,z_N)=(0,z_1,\ldots,z_N)\). The four coordinate lemmas establish these formulas directly; the subtraction defining \texttt{latentNoiseCLM} is an equality of continuous linear maps, so all later compositions are measurable.




\textbf{Correspondence with the displayed declarations.}


\begin{itemize}

\item \texttt{parameterNoiseCLM}: Retain the parameter coordinates and erase the latent coordinates.
\item \texttt{latentNoiseCLM}: Retain precisely the latent noise of the original particle scheme.
\end{itemize}


\noindent Lean declarations: \texttt{parameterNoiseCLM}, \texttt{parameterNoiseCLM\_parameter}, \texttt{parameterNoiseCLM\_particle}, \texttt{latentNoiseCLM}, \texttt{latentNoiseCLM\_parameter}, \texttt{latentNoiseCLM\_particle}.

\section{Compute the projection geometry and Gaussian moments}


The Euclidean square is \(\|P_\theta z\|^2=\sum_{j=1}^{d_\theta}z_j^2\), hence \(\|P_\theta z\|\le\|z\|\). A bounded linear map commutes with the Gaussian expectation, giving \(\mathbb E P_\theta\xi=0\). Each retained standard coordinate has second moment one; finite summation gives \(\mathbb E\|P_\theta\xi\|^2=d_\theta\).




The squared norm expands to \(\sum_{j=1}^{d_\theta}z_j^2\), with no latent terms. Every coordinate square is integrable because the standard Gaussian identity map belongs to \(L^2\). The coordinate variance identity gives \(\mathbb E z_j^2=1\) and the continuous-linear-map integration identity gives \(\mathbb E P_\theta Z=P_\theta\mathbb EZ=0\). Interchanging only a finite sum and the integral yields \(\mathbb E\|P_\theta Z\|^2=d_\theta\). No factor \(N\) belongs in this dimension count.



\noindent Lean declarations: \texttt{parameterNoise\_norm\_sq}, \texttt{parameterNoise\_lipschitz}, \texttt{integral\_parameterNoise\_stdGaussian}, \texttt{integral\_parameterNoise\_sq\_stdGaussian}.

\section{The product law cancels the mixed term exactly}


For any square-integrable \(g\) on a probability space independent of \(\xi\), expand \(\|g+bP_\theta\xi\|^2\). Fubini and bilinearity factor its mixed expectation as \(\langle\mathbb E g,\mathbb E P_\theta\xi\rangle=0\). Thus

\[\mathbb E\|g+bP_\theta\xi\|^2=\mathbb E\|g\|^2+b^2d_\theta.\]

The proof establishes integrability of every term before distributing the integral.




For a general \(g\in L^2(\mu)\), use the actual product law \(\mu\otimes\mathcal N(0,I)\), and expand \[\|g(x)+bP_\theta z\|^2=\|g(x)\|^2+2b\langle g(x),P_\theta z\rangle+b^2\|P_\theta z\|^2.\] The facts \texttt{hgsq} and \texttt{hzsq} prove integrability of the two squares on the product; \texttt{hcross} proves integrability of the bilinear term, so splitting the integral is justified. Independence here is the product-measure construction itself: \texttt{integral\_prod\_bilin} factors the mixed integral as \(\langle\mathbb Eg,\mathbb E P_\theta Z\rangle=0\). Hence \(\mathbb E\|g+bP_\theta Z\|^2=\mathbb E\|g\|^2+b^2d_\theta\), with no assumption that \(g\) is centered.




\textbf{Correspondence with the displayed declarations.}


\begin{itemize}

\item \texttt{integral\_add\_parameterNoise\_sq}: Exact independent-noise identity, for an arbitrary square-integrable source. The mixed term vanishes by the actual product Gaussian law.
\end{itemize}


\noindent Lean declarations: \texttt{integral\_add\_parameterNoise\_sq}.

\end{document}
